%% =====================================================================
%%  RESUME (2 pages) -- Research engineering / applied mathematics / industry
%%  Wöden Kusner
%% =====================================================================
%%  Build: pdflatex Kusner_Resume_Research-Engineering.tex (twice)
%%
%%  Design goals: one column, no tables or text boxes, plain section names,
%%  real text (not images), standard font, dates on the same line as the
%%  role -- so applicant-tracking systems parse it cleanly.
%%  Tailor the SUMMARY and the order of PROJECTS to each posting.
%% =====================================================================
\documentclass[11pt]{article}
\usepackage[letterpaper,margin=0.7in]{geometry}
\usepackage[expansion=false]{microtype}
\input{glyphtounicode}\pdfgentounicode=1
\usepackage{enumitem}
\usepackage{xcolor}
\usepackage[hidelinks]{hyperref}
\hypersetup{pdftitle={Woden Kusner - Resume}, pdfauthor={W\"oden Kusner (Woden Kusner)},
  pdfkeywords={Woden Kusner, applied mathematics, research engineer, C, Python, Julia,
  interval arithmetic, certified computation, optimization, computational geometry,
  computer-assisted proof, AI-assisted research}}
\pagestyle{empty}
\setlength{\parindent}{0pt}
\definecolor{rule}{gray}{0.55}
\hyphenpenalty=10000 \exhyphenpenalty=10000 \emergencystretch=3em  % no split words: cleaner ATS parsing

% --- helpers -----------------------------------------------------------
\newcommand{\sect}[1]{\par\vspace{7pt}{\large\bfseries #1}\par\vspace{1pt}%
  {\color{rule}\hrule height .5pt}\vspace{4pt}}
\newcommand{\role}[3]{\par\vspace{3pt}\textbf{#1}\ \hfill\mbox{#3}\par   % title, org, dates
  {\itshape #2}\par}
\setlist[itemize]{leftmargin=1.1em, itemsep=1pt, topsep=1pt, parsep=0pt, label=$\bullet$}

\begin{document}

%% ---------------- HEADER ----------------
{\LARGE\bfseries W\"oden Kusner}\par\vspace{2pt}
Applied mathematician $|$ certified computation, geometry, and optimization\par\vspace{2pt}
\href{mailto:wkusner@gmail.com}{wkusner@gmail.com} $|$
\href{https://wkusner.github.io}{wkusner.github.io} $|$
Appleton, WI $|$ U.S.\ citizen

%% ---------------- SUMMARY ----------------
\sect{Summary}
Mathematician (Ph.D., University of Pittsburgh, advisor Thomas C.~Hales) with twelve years of research in
discrete geometry, optimization, and computer-assisted proof. Builds rigorous numerical software
(interval arithmetic, branch-and-bound search, independent certificate checking) and turns hard technical
problems into verified results. Experienced explaining mathematics to non-specialists, mentoring, and
running a campus-wide quantitative support program.

%% ---------------- SKILLS ----------------
\sect{Technical Skills}
\begin{itemize}
  \item \textbf{Languages and tools:} Mathematica and Python (main tools; incl.\ mpmath); working familiarity
        with C (FLINT/Arb), Julia, and HTML/CSS/JavaScript; Git/GitHub; \LaTeX.
  \item \textbf{Methods:} interval arithmetic and certified computation (FLINT/Arb); branch-and-bound search;
        nonlinear optimization and critical-point analysis; computational geometry; discrepancy and
        quasi-Monte Carlo point sets; probability and statistics.
  \item \textbf{Workflow:} AI-assisted research and software development (Claude, Claude Code); reproducible
        releases with archived DOIs.
\end{itemize}

%% ---------------- PROJECTS ----------------
\sect{Selected Projects}
\begin{itemize}
  \item \textbf{kisscrit} (2026, open source). C library on FLINT/Arb for certified search of critical
        configurations of kissing spheres: interval branch-and-bound plus a separate certificate checker.
        Released v0.1.0 with DOI \href{https://doi.org/10.5281/zenodo.23603268}{10.5281/zenodo.23603268}.
  \item \textbf{Spherical cap discrepancy} (2 papers in preparation; code and data public). Exact optimal
        configurations of $N \le 6$ points on the sphere, with a computer-assisted proof for $N = 5, 6$
        cross-verified by two independent interval-arithmetic pipelines (Arb and mpmath.iv); open-source
        \texttt{capdisc} verification code, certificate checker, and an atlas of critical configurations
        for $N = 5$--$12$ (\href{https://github.com/wkusner/discrepancy}{github.com/wkusner/discrepancy}).
  \item \textbf{Critical radii of kissing configurations} (in preparation). Algebraicity and balance results
        for twelve-sphere configurations, certified with kisscrit.
  \item \textbf{Computer-assisted packing proofs.} Optimality of packings of regular pentagons, with T.~Hales
        (in revision, \emph{Trans.\ Amer.\ Math.\ Soc.}); local optimality of the double lattice packing,
        with Y.~Kallus (\emph{Discrete Comput.\ Geom.}, 2016); density bounds for cylinders and polycylinders
        (\emph{Discrete Comput.\ Geom.}, 2014, 2016).
  \item \textbf{Chirality measure for curves} (with G.~Dietler, E.~Rawdon, R.~Kusner, P.~Szymczak; in revision).
        Tensor-valued chirality measure built from hydrodynamic intuition; applies to polygonal curves,
        polymers, and biomolecules.
  \item \textbf{Hyperuniform point sets on the sphere} (with J.~Brauchart, P.~Grabner, J.~Ziefle;
        \emph{Constr.\ Approx.}\ 2019, \emph{Monatsh.\ Math.}\ 2020). Deterministic and random constructions of
        well-distributed point sets.
  \item \textbf{Of Instruments} (2026). Interactive browser versions of historical mathematical instruments
        (slide rules, planimeters, Napier's bones, verniers), written in JavaScript.
  \item \textbf{Community data infrastructure.} Fox Valley Data Exchange data-lakehouse specification working
        group (2025); Community Data Hub data team for Imagine Fox Cities (2021--2022).
\end{itemize}

%% ---------------- EXPERIENCE ----------------
\sect{Experience}
\role{Director, Quantitative Resource Center; Adjunct Assistant Professor of Mathematics}{Lawrence University}{2025--present}
\begin{itemize}
  \item Run quantitative peer tutoring for the college: recruit, train, and mentor peer tutors; coordinate
        support with departments; advise on ALEKS math placement.
  \item Created the Quantitative Learning Assistant program (2025): wrote the position description, hire and
        schedule assistants, and pair each with a course instructor.
  \item Provide individualized support for case-managed students and hold open office hours for students,
        faculty, and staff.
\end{itemize}
\role{Instructor of Mathematics}{Lawrence University}{2023--2025}
\begin{itemize}
  \item Designed the ALEKS-based course structure and exercise pack for Preparation for Calculus, later used by
        colleagues teaching the course; supervised directed studies and summer undergraduate research.
\end{itemize}
\role{Assistant Professor (non-tenure-track)}{University of Georgia}{2020--2022}
\begin{itemize}
  \item Taught Calculus I and II; UGA Teaching Academy Fellow; member of the instructional corps recognized by
        the 2022 Regents' Teaching Excellence Award.
\end{itemize}
\role{Assistant Professor (non-tenure-track) and Postdoctoral Associate}{Vanderbilt University}{2017--2020}
\begin{itemize}
  \item Research in the Center for Constructive Approximation; Data Science Institute faculty affiliate.
  \item Taught probability, statistical inference, and multivariable calculus; supervised teams of teaching
        assistants and graders; mentored undergraduate research.
  \item Organized the Computational Analysis Seminar; co-organized the Shanks Workshop on Computations and
        Linear Programming Bounds for Energy, Packing and Covering (2019).
\end{itemize}
\role{FWF Postdoctoral Researcher}{Graz University of Technology, Austria}{2014--2017}
\begin{itemize}
  \item Research on point configurations, discrepancy, and packing; taught a course on packings, lattices,
        and configurations.
\end{itemize}

%% ---------------- EDUCATION ----------------
\sect{Education}
\textbf{Ph.D., Mathematics}, University of Pittsburgh \hfill 2014\par
Dissertation: \emph{Bounds on packing density via slicing} (advisor Thomas C.~Hales)\par
\textbf{B.S., Mathematics}, Haverford College \hfill 2007

%% ---------------- RECOGNITION ----------------
\sect{Publications, Talks, and Professional Activity}
\begin{itemize}
  \item 7 peer-reviewed papers (incl.\ 3 in \emph{Discrete \& Computational Geometry}) and 6 papers in revision or preparation;
        full list at \href{https://wkusner.github.io}{wkusner.github.io}.
  \item 37 invited and contributed talks, including Oberwolfach, MSRI, ICERM, AIM, BIRS--CMO, and the
        Isaac Newton Institute.
  \item Referee for \emph{Annals of Mathematics}, \emph{Discrete \& Computational Geometry}, IMRN,
        \emph{SIAM J.\ Discrete Math.}, and others; NSF Graduate Research Fellowship reviewer (2022).
  \item Languages: English (native), German (intermediate), Spanish (elementary).
\end{itemize}

\end{document}
